\input{../preamble.tex}

\SetDate[19/03/2026]
\reversemarginpar
\setlength{\marginparsep}{.5cm}

\begin{document}
%\pagestyle{fancy}
%\fancyhead[L]{Seconde}
%\fancyhead[C]{\textbf{Évaluation blanche — Probabilités}}
%\fancyhead[R]{\today}

\marginpar{[$\infty$]}


\exe{, difficulty=1}{
	Remplir le tableau suivant de sorte que
		\begin{itemize}
			\item
			$P$ soit une loi de probabilité sur l'univers $\bigset{1 ; 2 ; 3 ; 4}$ ; et
			\item
			$P$ soit une fonction affine sur $\R$ non constante (sinon c'est trop facile).
		\end{itemize}
	\begin{center}
	\begin{tabular}{|c|c|c|c|c|} \hline
		$x$ & 1 & 2 & 3 & 4 \\ \hline
		$P(x)$ & & & & \\ \hline
	\end{tabular}
	\end{center}
	Il existe une infinité de résultats possibles.
}{exe:3}{
%	Il s'agit de trouver $a$ et $b$ tels que $P(x) = ax+b$ soit une loi de probabilité sur $\bigset{1 ; 2 ; 3 ;4}$.
%	On a le tableau de valeurs suivant.
%	\begin{center}
%	\begin{tabular}{|c|c|c|c|c|} \hline
%		$x$ & 1 & 2 & 3 & 4 \\ \hline
%		$P(x)$ & $a+b$ & $2a+b$ & $3a+b$ & $4a+b$ \\ \hline
%	\end{tabular}
%	\end{center}
%	Il faut que chaque probabilité soit entre 0 et 1 et que la somme soit égale à 1.
%	
%	La somme égale à 1 implique que
%		\begin{align*}
%			(a+b)+(2a+b)+(3a+b)+(4a+b) &= 1 \\
%			10a + 4b &= 1 \\
%			b &= \frac14 - \frac52 a
%		\end{align*}
%	Remarquons que $b$ est une fonction affine en $a$ !
%	Il existe donc une infinité de tels paramètres $(a ; b)$, chacun un point appartenant à la droite $y = -\frac52 x + \frac14$.
%	
%	Imposons désormais que chaque probabilité appartient à $[0 ; 1]$.
%	La première probabilité, $a+b$ contraint $a$ comme suit.
%		\begin{alignat*}{2}
%			0 &\leq a+b &\leq 1 \\
%			0 &\leq a + \frac14 - \frac52 a &\leq 1 \\
%			0 &\leq \frac14 - \frac32 a &\leq 1 \\
%			-\frac14 &\leq -\frac32 a &\leq \frac34 \\
%			\frac16 &\geq a &\geq -\frac12 \\
%		\end{alignat*}
%	La deuxième probabilité, $4a+b$ contraint $a$ comme suit.
%		\begin{alignat*}{2}
%			0 &\leq 4a+b &\leq 1 \\
%			0 &\leq 4a + \frac14 - \frac52 a &\leq 1 \\
%			0 &\leq \frac14 + \frac32 a &\leq 1 \\
%			-\frac14 &\leq \frac32 a &\leq \frac34 \\
%			-\frac16 &\leq a &\leq \frac12
%		\end{alignat*}
%	Comme pour passer d'une probabilité à l'autre, on ajoute $a$, on peut se convaincre sans problème que si $a+b \in [0 ;1]$ et $4a+b \in [0;1]$, alors chaque probabilité appartient aussi à $[0 ; 1]$.
%	
%	En conclusion, il faut donc que 
%		\[ a \in \left[-\frac12 ; \frac16\right] \cap \left[-\frac16 ; \frac12\right] = \left[-\frac16 ; \frac16\right]. \]
%	Prenons $a = \frac1{25}$ car on aime les fractions, ce qui implique
%		\[ b = \frac14 - \frac52 a = \frac14 - \frac1{10} = \frac3{20}. \]
%	Par conséquent, $P(x) = \frac1{25}x + \frac3{20} = \frac{4x + 15}{100}$ de tableau de valeurs
%	\begin{center}
%	\begin{tabular}{|c|c|c|c|c|} \hline
%		$x$ & 1 & 2 & 3 & 4 \\ \hline
%		$P(x)$ & $\frac{19}{100}$ & $\frac{23}{100}$ & $\frac{27}{100}$ & $\frac{31}{100}$ \\ \hline
%	\end{tabular}
%	\end{center}
%	Ouf !
%
%	\underline{Deuxième version}
	
	Avec un tableau de valeurs à valeurs toutes positives, il suffit de diviser par la somme des images pour les normaliser : elles deviendront toutes inférieures à 1, et leur somme vaudra 1.
	C'est d'ailleurs le principe du calcul de fréquences.
	
	On peut donc poser $a=1$ et $b=0$ de telle sorte à avoir uniquement des valeurs positives dans notre tableau, puis diviser par la somme $1+2+3+4 = 10$.
	Ainsi $P(x) = \frac{x}{10}$ fonctionne !
	\begin{center}
	\begin{tabular}{|c|c|c|c|c|} \hline
		$x$ & 1 & 2 & 3 & 4 \\ \hline
		$P(x)$ & $\frac{1}{10}$ & $\frac{2}{10}$ & $\frac{3}{10}$ & $\frac{4}{10}$ \\ \hline
	\end{tabular}
	\end{center}
	
	En prenant $a = 1$ et $b=34$, la somme des images vaut $35+36+37+38=146$.
	Donc $P(x) = \frac{x+34}{146}$ fonctionne.
	
	En général,
		\[ P(x) = \frac{f(x)}{f(1)+f(2)+f(3)+f(4)} \]
	fonctionne tant que $f(x) \geq 0$ pour tout $x\in\{1;2;3;4\}$ (et que le dénominateur n'est pas nul).
}

%%%%%%%%%%%

\label{lastpage}
\newpage
\fancyhead[C]{\textbf{Solutions}}
\shipoutAnswer
	
\end{document}
